% ============================================================
%  GUÍA 4 — TRIGONOMETRÍA APLICADA
%  Basada en la plantilla institucional Usach Premium
% ============================================================

\documentclass[letterpaper, 12pt]{article}

% ============================================================
% PAQUETES
% ============================================================
\usepackage[utf8]{inputenc}
\usepackage[spanish]{babel}
\usepackage{amsmath, amssymb, amsfonts}
\usepackage{geometry}
\usepackage{fancyhdr}
\usepackage{enumitem}
\usepackage{graphicx}
\usepackage{hyperref}
\usepackage{multicol}
\usepackage{parskip}
\usepackage{xcolor}
\usepackage{tcolorbox}
\usepackage{eso-pic}
\usepackage{transparent}
\usepackage{tikz}
\usetikzlibrary{patterns, shadows.blur, fadings, calc, arrows.meta, decorations.pathreplacing}
\usepackage{array}
\usepackage{colortbl}
\tcbuselibrary{skins, shadows.blur}

% ============================================================
% GEOMETRÍA
% ============================================================
\geometry{letterpaper, margin=1.5cm, top=3cm, bottom=2.5cm, headheight=2cm}

% ============================================================
% COLORES INSTITUCIONALES
% ============================================================
\definecolor{celesteSuave}{HTML}{F0F7FF}
\definecolor{azulFuerte}{HTML}{1A5276}
\definecolor{turquesaUsach}{HTML}{33CCCC}
\definecolor{enlace}{HTML}{1A5276}
% Colores adicionales usados en figuras
\definecolor{azulMedio}{HTML}{2980B9}
\definecolor{marronMadera}{HTML}{935116}
\definecolor{verdePro}{HTML}{1E8449}

% ============================================================
% RUTAS DE IMÁGENES
% ============================================================
\newcommand{\logoup}{./logo_up.png}
\newcommand{\logousach}{./logo_usach.png}
\newcommand{\logofondo}{./logo_up.png}

% ============================================================
% INFORMACIÓN DE LA GUÍA
% ============================================================
\newcommand{\institucion}{Usach Premium}
\newcommand{\curso}{Cálculo I}
\newcommand{\guiasnumero}{Guía 4 - Trigonometría Aplicada}
\newcommand{\semestre}{1er. Semestre de 2026}
\newcommand{\preparadopor}{Martina Copia — Usach Premium.}
\newcommand{\webinstitucion}{www.usachpremium.cl}
\newcommand{\instagraminstitucion}{https://www.instagram.com/usach.premium}
\newcommand{\transparencia}{0.05}
\newcommand{\fraseportada}{"Por y para estudiantes"}

\newcommand{\listacontenidos}{
    \item[\color{azulFuerte}$\Rightarrow$] Razones trigonométricas
    \item[\color{azulFuerte}$\Rightarrow$] Identidades de suma y resta
    \item[\color{azulFuerte}$\Rightarrow$] Tabla de valores notables
    \item[\color{azulFuerte}$\Rightarrow$] Identidades trigonométricas
}

% ============================================================
% HYPERREF
% ============================================================
\hypersetup{
    colorlinks=true,
    linkcolor=enlace,
    urlcolor=enlace,
    citecolor=enlace,
    pdfborder={0 0 0},
}

% ============================================================
% ENCABEZADO Y PIE DE PÁGINA
% ============================================================
\pagestyle{fancy}
\fancyhf{}
\fancyhead[L]{%
    \includegraphics[height=1.2cm]{\logoup}%
}
\fancyhead[R]{%
    \begin{minipage}[b]{0.42\textwidth}
        \raggedleft
        \fontsize{9pt}{11pt}\selectfont
        \textbf{\color{azulFuerte}\institucion} \\
        \curso\ — \guiasnumero \\
        \textit{\color{gray}@usach.premium}
    \end{minipage}\hspace{0.1cm}%
    \includegraphics[height=1.2cm]{\logousach}%
}
\fancyfoot[L]{\fontsize{9pt}{11pt}\selectfont \textit{\fraseportada}}
\fancyfoot[R]{\fontsize{9pt}{11pt}\selectfont Página \thepage}
\renewcommand{\headrulewidth}{0.4pt}
\renewcommand{\footrulewidth}{0.4pt}

\fancypagestyle{plain}{
    \fancyhf{}
    \renewcommand{\headrulewidth}{0pt}
    \renewcommand{\footrulewidth}{0pt}
}

% ============================================================
% MARCA DE AGUA
% ============================================================
\newcommand{\MarcaAgua}{
    \AddToShipoutPicture{
        \AtPageCenter{
            \makebox(0,0){%
                \transparent{\transparencia}%
                \includegraphics[width=0.7\textwidth]{\logofondo}%
            }
        }
    }
}

% ============================================================
% CAJAS REUTILIZABLES
% ============================================================

\newenvironment{kitpropiedades}[1]{%
    \begin{tcolorbox}[colback=celesteSuave, colframe=azulFuerte,
        title=\textbf{#1}, fonttitle=\bfseries]
}{%
    \end{tcolorbox}
}

\newtcolorbox{desafiobox}[1]{
    colback=white,
    colframe=azulFuerte,
    fonttitle=\bfseries,
    coltitle=white,
    title=#1,
    arc=8pt,
    boxrule=1.2pt,
    enhanced,
    drop shadow={black!5}
}

\newtcolorbox{ejerciciobox}[1]{
    enhanced, boxrule=1pt, colframe=azulFuerte, colback=white,
    sharp corners, drop shadow={black!10!white},
    title=\textbf{#1}, coltitle=white, colbacktitle=azulFuerte,
    fonttitle=\small\sffamily, attach boxed title to top left={xshift=0.3cm, yshift=-3mm},
    boxed title style={sharp corners, boxrule=0.5pt, colframe=azulFuerte}
}

% ============================================================
% ENTORNOS DE SECCIÓN
% ============================================================

\newenvironment{introduccion}{%
    \section*{Introducción}%
    \MarcaAgua
}{}

\newenvironment{ejercicios}{%
    \MarcaAgua
}{}

% ============================================================
\begin{document}
\thispagestyle{plain}

% ============================================================
% PORTADA
% ============================================================
\AddToShipoutPicture*{%
    \put(54, 700){\includegraphics[width=2cm]{\logoup}}
    \put(420, 706){\includegraphics[width=5cm]{\logousach}}
}

\vspace*{0.5cm}

\begin{center}
    \color{azulFuerte}
    {\huge \textbf{GUÍA DE EJERCICIOS}} \\[0.5cm]
    {\Huge \textbf{\curso}} \\[0.3cm]
    {\Large \guiasnumero} \\[1.5cm]

    \begin{tcolorbox}[colback=celesteSuave, colframe=azulFuerte, arc=10pt,
        width=0.85\textwidth, center, boxrule=1.5pt]
        \centering
        \vspace{0.2cm}
        {\Large \textbf{Contenidos}}
        \vspace{0.3cm}
        \color{black}
        \begin{itemize}[leftmargin=3cm]
            \listacontenidos
        \end{itemize}
        \vspace{0.2cm}
    \end{tcolorbox}
\end{center}

\vspace{1.5cm}

\begin{center}
    {\Large \textbf{\semestre}} \\[0.8cm]
    {\large \textbf{\preparadopor}}
\end{center}

\vspace{2cm}

\begin{center}
    {\color{azulFuerte}\hrule height 0.5pt}
    \vspace{0.3cm}
    \begin{minipage}{0.45\textwidth}
        \centering
        \textbf{Instagram:} \\
        \small\href{\instagraminstitucion}{@usach.premium}
    \end{minipage}
    \begin{minipage}{0.45\textwidth}
        \centering
        \textbf{Web:} \\
        \small\href{http://\webinstitucion}{\webinstitucion}
    \end{minipage}
\end{center}

\pagenumbering{arabic}
\newpage

% ============================================================
% KIT DE PROPIEDADES
% ============================================================
\begin{ejercicios}

\begin{kitpropiedades}{RESUMEN: FÓRMULAS Y VALORES NOTABLES}
\begin{multicols}{2}
    \textbf{Definiciones de Razones Trigonométricas:}
    \begin{itemize}[leftmargin=10pt]
        \item $\sen \alpha = \frac{op}{hip} = \frac{a}{c} \quad \csc \alpha = \frac{hip}{op} = \frac{c}{a}$
        \item $\cos \alpha = \frac{ady}{hip} = \frac{b}{c} \quad \sec \alpha = \frac{hip}{ady} = \frac{c}{b}$
        \item $\tan \alpha = \frac{op}{ady} = \frac{a}{b} \quad \cot \alpha = \frac{ady}{op} = \frac{b}{a}$
    \end{itemize}

    \columnbreak

    \textbf{Identidades de Suma y Resta:}
    \begin{itemize}[leftmargin=10pt]
        \item $\sen(\alpha \pm \beta) = \sen\alpha \cos\beta \pm \sen\beta \cos\alpha$
        \item $\cos(\alpha \pm \beta) = \cos\alpha \cos\beta \mp \sen\alpha \sen\beta$
    \end{itemize}
\end{multicols}

\vspace{0.2cm}

\begin{center}
    \small
    \renewcommand{\arraystretch}{1.2}
    \begin{tabular}{|c|c|c|c|c|c|}
    \hline
    \rowcolor{azulFuerte} \color{white} $\alpha$ & \color{white} $0^\circ$ & \color{white} $30^\circ$ & \color{white} $45^\circ$ & \color{white} $60^\circ$ & \color{white} $90^\circ$ \\ \hline
    $\sen(\alpha)$ & $0$ & $1/2$ & $\sqrt{2}/2$ & $\sqrt{3}/2$ & $1$ \\ \hline
    $\cos(\alpha)$ & $1$ & $\sqrt{3}/2$ & $\sqrt{2}/2$ & $1/2$ & $0$ \\ \hline
    $\tan(\alpha)$ & $0$ & $\sqrt{3}/3$ & $1$ & $\sqrt{3}$ & $\not\exists$ \\ \hline
    $\csc(\alpha)$ & $\not\exists$ & $2$ & $\sqrt{2}$ & $2\sqrt{3}/3$ & $1$ \\ \hline
    $\sec(\alpha)$ & $1$ & $2\sqrt{3}/3$ & $\sqrt{2}$ & $2$ & $\not\exists$ \\ \hline
    $\cot(\alpha)$ & $\not\exists$ & $\sqrt{3}$ & $1$ & $\sqrt{3}/3$ & $0$ \\ \hline
    \end{tabular}
\end{center}
\end{kitpropiedades}

\vspace{0.4cm}

% ============================================================
% BLOQUE DE PROBLEMAS
% ============================================================

{\small\color{gray}\textit{Nota: algunos ejercicios de esta guía fueron extraídos del material oficial de la Coordinación de Cálculo I. Se dan los créditos correspondientes a sus autores.}}

\vspace{0.3cm}

\begin{ejerciciobox}{Ejercicio 1}
\centering
\begin{minipage}[c]{0.58\textwidth}
    1. Calcula la altura de un árbol que a una distancia de 10 m se ve bajo un ángulo de elevación de 30°.
\end{minipage}
\begin{minipage}[c]{0.38\textwidth}
    \centering
    \begin{tikzpicture}[scale=0.7, >=Stealth] 
        % Árbol a la derecha
        \shade[left color=marronMadera!90!black, right color=marronMadera!60] (3.8,0) rectangle (4,2.2);
        \foreach \p in {(3.6,2.4), (3.9,3), (4.2,2.4), (3.9,2.3)} {\shade[ball color=verdePro!80] \p circle (0.55);}
        % Triángulo hacia la izquierda
        \draw[azulMedio, thick, dashed] (0,0) -- (3.9,2.4) -- (3.9,0) -- cycle;
        \draw[azulMedio, thick] (0.6,0) arc (0:30:0.6);
        \node[azulFuerte] at (1,0.3) {\tiny 30°}; 
        \draw[|<->|, azulFuerte] (0,-0.4) -- (3.9,-0.4) node[midway, fill=white, font=\tiny] {10 m};
    \end{tikzpicture}
\end{minipage}
\end{ejerciciobox}

\begin{ejerciciobox}{Ejercicio 2}
\centering
\begin{minipage}[c]{0.52\textwidth}
    2. Calcule la altura $h_{c}$ del siguiente triángulo equilátero de lado 2.
\end{minipage}
\begin{minipage}[c]{0.44\textwidth}
    \centering
    \begin{tikzpicture}[scale=1.5, >=Stealth]
        \fill[yellow!20] (0,0) -- (1,1.73) -- (1,0) -- cycle;
        \fill[blue!15] (1,0) -- (1,1.73) -- (2,0) -- cycle;
        \draw[thick] (0,0) node[below left] {\small A} -- (2,0) node[below right] {\small B} -- (1,1.73) node[above] {\small C} -- cycle;
        \draw[dashed] (1,1.73) -- (1,0) node[midway, right] {\small $h_c$};
        \draw (0.3,0) arc (0:60:0.3); \node at (0.2,0.15) {\tiny 60°};
        \draw (1.7,0) arc (180:120:0.3); \node at (1.8,0.15) {\tiny 60°};
        \draw (0.85,1.47) arc (240:300:0.3);
        \node at (0.88,1.35) {\tiny $x$}; \node at (1.12,1.35) {\tiny $x$};
        \draw (0.85,0) -- (0.85,0.15) -- (1,0.15);
        \node[left] at (0.5,0.9) {\small 2}; \node[right] at (1.5,0.9) {\small 2};
        \draw [thick, decoration={brace, mirror, raise=5pt}, decorate] (0,0) -- (1,0);
        \draw [thick, decoration={brace, mirror, raise=5pt}, decorate] (1,0) -- (2,0);
    \end{tikzpicture}
\end{minipage}
\end{ejerciciobox}

\begin{ejerciciobox}{Ejercicio 3}
\centering
\begin{minipage}[c]{0.58\textwidth}
    3. Para medir la altura de un cerro se realizan dos mediciones. Desde un punto A el ángulo es de 45°. Luego, 0,5 km más lejos, el ángulo es de 30°. Determine la altura $h$.
\end{minipage}
\begin{minipage}[c]{0.38\textwidth}
    \centering
    \begin{tikzpicture}[scale=0.55, >=Stealth] 
        % Cerro a la izquierda
        \shade[top color=gray!10, bottom color=gray!50] (0,0) -- (1.3,3) -- (2.5,0);
        \draw[azulFuerte, thick] (0,0) -- (1.3,3) -- (2.5,0);
        \draw[thick] (0,0) -- (6.5,0);
        
        % Línea de altura h
        \draw[dashed, gray] (1.3,3) -- (1.3,0);
        \node[right, azulFuerte] at (1.3,1.5) {\tiny $h$};

        % Mediciones hacia la derecha
        \draw[azulMedio, dashed] (6.5,0) -- (1.3,3); % Ángulo 30
        \draw[azulMedio, dashed] (3.7,0) -- (1.3,3); % Ángulo 45
        
        % Ángulos
        \draw (5.7,0) arc (180:150:0.8); \node at (5.3,0.4) {\tiny 30°};
        \draw (3.2,0) arc (180:129:0.5); \node at (2.9,0.5) {\tiny 45°};
        
        % Cotas
        \draw[|<->|, gray] (3.7,-0.6) -- (6.5,-0.6) node[midway, fill=white, font=\tiny] {0,5 km};
        \draw[|<->|, azulMedio] (1.3,-0.6) -- (3.7,-0.6) node[midway, fill=white, font=\tiny] {$x$};
    \end{tikzpicture}
\end{minipage}
\end{ejerciciobox}

\begin{ejerciciobox}{Ejercicio 4}
\centering
\begin{minipage}[c]{0.58\textwidth}
    4. Un carpintero corta el borde de un tablero de 3'' de largo, con una inclinación de 30° de la vertical, empezando desde un punto ubicado a 3/4'' del borde. Determine la longitud del corte diagonal.
\end{minipage}
\begin{minipage}[c]{0.38\textwidth}
    \centering
    \begin{tikzpicture}[scale=0.8, >=Stealth]
        \fill[marronMadera!15] (0,0) rectangle (3,2.5);
        \draw[azulFuerte, thick] (0,0) rectangle (3,2.5);
        \draw[red!90!black, ultra thick] (0.8,2.5) -- (2.3,0) node[midway, right, xshift=2pt] {\small $y$};
        \draw[dashed, gray] (0.8,2.5) -- (0.8,0);
        \draw (0.8,1.9) arc (-90:-61:0.6); \node at (1.1,1.7) {\tiny 30°};
        \draw[|<->|, gray] (0.8,-0.4) -- (2.3,-0.4) node[midway, below] {\small $x$};
        \draw[|<->|, azulFuerte] (0,2.8) -- (0.8,2.8) node[midway, above, font=\tiny] {3/4''};
        \draw[|<->|, azulFuerte] (-0.3,0) -- (-0.3,2.5) node[midway, rotate=90, above, font=\tiny] {3''};
    \end{tikzpicture}
\end{minipage}
\end{ejerciciobox}

\vspace{0.4cm}
\begin{center}
    \color{azulFuerte}\textbf{\small IDENTIDADES TRIGONOMÉTRICAS}
    \hrule height 0.5pt
\end{center}

\begin{ejerciciobox}{Ejercicios 5 y 6}
    \vspace{0.3cm}
    Resuelva las siguientes identidades trigonométricas:
    \begin{enumerate}[label=\textbf{\alph*)}]
        \item $\tan(\theta) \cdot \cot(\theta) - \cos^{2}(\theta) = \sen^{2}(\theta)$
        \item $\sen(\theta)(\cot(\theta) + \tan(\theta)) = \sec(\theta)$
    \end{enumerate}
\end{ejerciciobox}

\end{ejercicios}

\end{document}
















